% !TeX program = xelatex
% 华数杯数学建模竞赛模板 — 基于 JXUSTmodeling 文档类
\documentclass{JXUSTmodeling}
% ⛔ booktabs, multirow, graphicx, caption, subcaption, placeins 已由 cls 加载，不要重复加载
\usepackage{tikz}
\usetikzlibrary{arrows.meta,positioning,shapes.geometric,fit,calc}
\usepackage{algorithm2e}
\usepackage{float}
\usepackage{needspace}
\usepackage{longtable}

% === 引用上标格式 ===
\makeatletter
\def\@cite#1#2{\textsuperscript{[{#1\if@tempswa , #2\fi}]}}
\makeatother

\biaoti{[论文标题]}
\keyword{[关键词1]、[关键词2]、[关键词3]}

\begin{document}

% === 华数杯报名信息表（封面，电子版匿名提交：所属类别 / 参赛编号 留空）===
{\centering
\begin{tabular}{|>{\centering\arraybackslash}p{0.14\textwidth}|>{\centering\arraybackslash}p{0.58\textwidth}|>{\centering\arraybackslash}p{0.14\textwidth}|}
\hline
\zihao{5}所属类别 & \multirow{3}{0.58\textwidth}{\centering\zihao{4}\bfseries 2026年“华数杯”全国大学生数学建模竞赛} & \zihao{5}参赛编号 \\
\cline{1-1}\cline{3-3}
 & & \\
 & & \\
\hline
\end{tabular}
\par}

\vskip2ex

% === 摘要（cls 自动生成标题和关键词）===
\begin{abstract}

[中文摘要内容：问题概述 + 每个子问题的方法和数值结果 + 结论]

\end{abstract}

% === 正文 ===
\input{sections/1_restatement}
\input{sections/2_analysis}
\input{sections/3_assumptions}
\input{sections/4_symbols}
\input{sections/5_problem1}
\input{sections/6_problem2}
\input{sections/7_problem3}
\input{sections/8_evaluation}

% === 参考文献 ===
\begin{thebibliography}{99}
% \bibitem{ref1} 作者. 标题[J]. 期刊, 年份.
\end{thebibliography}

% === 附录 ===
\begin{appendixx}
\input{sections/A_code}
\end{appendixx}

\end{document}
